% ============================================================
\chapter{Discrete Random Variables}
\label{ch:discrete}
% ============================================================

Rolling a die, counting customers in a queue, observing radioactive decays per second: in each of these situations, chance produces an outcome that can be enumerated. Discrete random variables---those taking only finitely or countably many values---are the first we encounter, and the most intuitive. Yet they hold surprises: the Poisson distribution, discovered by Sim\'eon Denis Poisson in 1837 to model judicial errors, reappears everywhere, from call centres to V-2 bomb impacts on London. The geometric distribution captures the patience of the stubborn gambler, while the binomial distribution, direct heir to Jacob Bernoulli's work, remains the reference model for repeated experiments.

This chapter builds the formal framework for discrete random variables and develops the classical distributions that populate all of probability theory.

\begin{intuition}
A random variable is a function that assigns a numerical value to the outcome of
a random experiment. This chapter treats the discrete case: the set of values
taken is finite or countable.
\end{intuition}

% ------------------------------------------------------------
\section{Definition and Distribution}
% ------------------------------------------------------------

\begin{definition}[Random variable]
Let $(\Omega, \mathcal{F}, \PP)$ be a probability space. A \textbf{real-valued
random variable} is a measurable function $X : \Omega \to \R$, i.e.\ for every
$B \in \mathcal{B}(\R)$:
\[
    X^{-1}(B) = \{\omega \in \Omega : X(\omega) \in B\} \in \mathcal{F}.
\]
\end{definition}

\begin{definition}[Discrete random variable]
A random variable $X$ is \textbf{discrete} if it takes values in a finite or
countable set $\{x_1, x_2, \ldots\} \subseteq \R$.
\end{definition}

\begin{definition}[Probability mass function]
The \textbf{distribution} of $X$ is the probability measure $P_X$ on
$(\R, \mathcal{B}(\R))$ defined by $P_X(B) = \PP(X \in B)$. For a discrete
variable, the distribution is completely determined by the \textbf{probability
mass function} (pmf):
\[
    p_X(x_k) = \PP(X = x_k), \quad k = 1, 2, \ldots
\]
with $p_X(x_k) \geq 0$ and $\sum_k p_X(x_k) = 1$.
\end{definition}

\begin{definition}[Cumulative distribution function]
The \textbf{cumulative distribution function} (cdf) of $X$ is:
\[
    F_X(x) = \PP(X \leq x) = \sum_{x_k \leq x} p_X(x_k), \quad x \in \R.
\]
\end{definition}

\begin{proposition}[Properties of $F_X$]
\begin{enumerate}[label=(\roman*)]
    \item $F_X$ is non-decreasing.
    \item $\lim_{x \to -\infty} F_X(x) = 0$ and $\lim_{x \to +\infty} F_X(x) = 1$.
    \item $F_X$ is right-continuous.
    \item $\PP(X = a) = F_X(a) - F_X(a^-)$ (the jump at $a$).
    \item For a discrete variable, $F_X$ is a step function.
\end{enumerate}
\end{proposition}

% ------------------------------------------------------------
\section{Classical Discrete Distributions}
% ------------------------------------------------------------

\subsection{Bernoulli distribution}

\begin{definition}[Bernoulli]
$X \sim \mathrm{Bernoulli}(p)$ with $p \in [0,1]$:
\[
    \PP(X = 1) = p, \quad \PP(X = 0) = 1 - p = q.
\]
$\E[X] = p$, $\Var(X) = pq$.
\end{definition}

\subsection{Binomial distribution}

\begin{definition}[Binomial]
$X \sim \mathrm{Bin}(n, p)$: the number of successes in $n$ independent
Bernoulli trials.
\[
    \PP(X = k) = \binom{n}{k} p^k (1-p)^{n-k}, \quad k = 0, 1, \ldots, n.
\]
$\E[X] = np$, $\Var(X) = np(1-p)$.
\end{definition}

\begin{example}
A fair coin is tossed 10 times. The probability of exactly 6 heads is:
\[
    \PP(X = 6) = \binom{10}{6}\left(\frac{1}{2}\right)^{10}
    = \frac{210}{1024} \approx 0.205.
\]
\end{example}

\subsection{Geometric distribution}

\begin{definition}[Geometric]
$X \sim \mathrm{Geom}(p)$: the number of trials until the first success.
\[
    \PP(X = k) = (1-p)^{k-1} p, \quad k = 1, 2, 3, \ldots
\]
$\E[X] = 1/p$, $\Var(X) = (1-p)/p^2$.
\end{definition}

\begin{proposition}[Memorylessness]
The geometric distribution is the only discrete distribution on $\N^*$ satisfying
the \textbf{memoryless property}:
\[
    \PP(X > m + n \mid X > m) = \PP(X > n) \quad \forall\, m, n \geq 0.
\]
\end{proposition}

\begin{proof}
$\PP(X > k) = (1-p)^k$. Hence:
\[
    \PP(X > m+n \mid X > m) = \frac{(1-p)^{m+n}}{(1-p)^m} = (1-p)^n = \PP(X > n). \qedhere
\]
\end{proof}

\subsection{Poisson distribution}

\begin{definition}[Poisson]
$X \sim \mathrm{Poisson}(\lambda)$ with $\lambda > 0$:
\[
    \PP(X = k) = e^{-\lambda}\frac{\lambda^k}{k!}, \quad k = 0, 1, 2, \ldots
\]
$\E[X] = \lambda$, $\Var(X) = \lambda$.
\end{definition}

\begin{theorem}[Poisson approximation]
If $X_n \sim \mathrm{Bin}(n, p_n)$ with $np_n \to \lambda$ as $n \to \infty$,
then for every $k \in \N$:
\[
    \PP(X_n = k) \xrightarrow[n \to \infty]{} e^{-\lambda}\frac{\lambda^k}{k!}.
\]
\end{theorem}

\begin{proof}
\begin{align*}
    \PP(X_n = k) &= \binom{n}{k} p_n^k (1-p_n)^{n-k} \\
    &= \frac{n!}{k!(n-k)!}\left(\frac{\lambda}{n}\right)^k
       \left(1 - \frac{\lambda}{n}\right)^{n-k} \\
    &= \frac{\lambda^k}{k!}\cdot
       \underbrace{\frac{n(n-1)\cdots(n-k+1)}{n^k}}_{\to 1}\cdot
       \underbrace{\left(1 - \frac{\lambda}{n}\right)^n}_{\to e^{-\lambda}}\cdot
       \underbrace{\left(1 - \frac{\lambda}{n}\right)^{-k}}_{\to 1}
    \xrightarrow[n\to\infty]{} e^{-\lambda}\frac{\lambda^k}{k!}. \qedhere
\end{align*}
\end{proof}

\begin{intuition}[title=\textbf{When to use Poisson?}]
The Poisson distribution models the number of rare events in a given interval
(phone calls, radioactive particles, failures, etc.). Rule of thumb: if $n \geq 30$
and $p \leq 0.1$, approximate $\mathrm{Bin}(n,p)$ by $\mathrm{Poisson}(np)$.
\end{intuition}

\subsection{Hypergeometric distribution}

\begin{definition}[Hypergeometric]
Draw $n$ items without replacement from a population of $N$ items, $K$ of which
are marked. $X$ = number of marked items among the $n$ drawn:
\[
    \PP(X = k) = \frac{\binom{K}{k}\binom{N-K}{n-k}}{\binom{N}{n}},
    \quad k = \max(0, n+K-N), \ldots, \min(n, K).
\]
$\E[X] = nK/N$.
\end{definition}

% ------------------------------------------------------------
\section{Summary Table}
% ------------------------------------------------------------

\begin{keyformulas}[title=\textbf{Fundamental discrete distributions}]
\begin{center}
\renewcommand{\arraystretch}{1.4}
\begin{tabular}{lccc}
    \toprule
    \textbf{Distribution} & \textbf{Parameters} & $\E[X]$ & $\Var(X)$ \\
    \midrule
    Bernoulli$(p)$ & $p \in [0,1]$ & $p$ & $p(1-p)$ \\
    Binomial$(n,p)$ & $n \in \N^*,\; p \in [0,1]$ & $np$ & $np(1-p)$ \\
    Geometric$(p)$ & $p \in (0,1]$ & $1/p$ & $(1-p)/p^2$ \\
    Poisson$(\lambda)$ & $\lambda > 0$ & $\lambda$ & $\lambda$ \\
    Uniform$\{1,\ldots,n\}$ & $n \in \N^*$ & $(n+1)/2$ & $(n^2-1)/12$ \\
    \bottomrule
\end{tabular}
\end{center}
\end{keyformulas}

% ------------------------------------------------------------
\section{Probability Generating Functions}
% ------------------------------------------------------------

\begin{definition}[Probability generating function]
Let $X$ be a random variable taking values in $\N$. The \textbf{probability
generating function} (pgf) of $X$ is:
\[
    G_X(s) = \E[s^X] = \sum_{k=0}^{\infty} \PP(X = k)\, s^k, \quad \abs{s} \leq 1.
\]
\end{definition}

\begin{proposition}[Properties of $G_X$]
\begin{enumerate}[label=(\roman*)]
    \item $G_X(1) = 1$.
    \item $\E[X] = G_X'(1)$.
    \item $\E[X(X-1)] = G_X''(1)$, whence
          $\Var(X) = G_X''(1) + G_X'(1) - [G_X'(1)]^2$.
    \item $\PP(X = k) = G_X^{(k)}(0)/k!$.
    \item If $X$ and $Y$ are independent, $G_{X+Y}(s) = G_X(s)\cdot G_Y(s)$.
\end{enumerate}
\end{proposition}

\begin{example}[Pgf of the Poisson distribution]
If $X \sim \mathrm{Poisson}(\lambda)$:
\[
    G_X(s) = \sum_{k=0}^{\infty} e^{-\lambda}\frac{\lambda^k}{k!}\,s^k
    = e^{-\lambda}\sum_{k=0}^{\infty}\frac{(\lambda s)^k}{k!}
    = e^{-\lambda}\cdot e^{\lambda s} = e^{\lambda(s-1)}.
\]
We recover $G_X'(1) = \lambda = \E[X]$ and $G_X''(1) = \lambda^2$, giving
$\Var(X) = \lambda^2 + \lambda - \lambda^2 = \lambda$.
\end{example}

% ------------------------------------------------------------
\section{Exercises}
% ------------------------------------------------------------

\begin{exercise}
Let $X \sim \mathrm{Bin}(n,p)$. Show directly (by calculation) that
$\E[X] = np$ and $\Var(X) = np(1-p)$.
\end{exercise}

\begin{exercise}
Let $X \sim \mathrm{Geom}(p)$. Compute $\E[X]$ using the probability
generating function.
\end{exercise}

\begin{exercise}
Let $X \sim \mathrm{Poisson}(\lambda)$ and $Y \sim \mathrm{Poisson}(\mu)$,
independent. Show that $X + Y \sim \mathrm{Poisson}(\lambda + \mu)$ using
probability generating functions.
\end{exercise}

\begin{exercise}
A telephone exchange receives an average of 4 calls per minute, modelled by a
Poisson distribution. Find the probability of:
\begin{enumerate}[label=(\alph*)]
    \item exactly 3 calls in one minute;
    \item at least 1 call in 30 seconds;
    \item more than 10 calls in 2 minutes.
\end{enumerate}
\end{exercise}

\begin{exercise}
A bag contains 5 red and 15 blue balls. Four balls are drawn without replacement.
Find the probability of getting exactly 2 red balls.
\end{exercise}

\begin{exercise}
A fair die is rolled until a 6 appears. Let $X$ be the number of rolls.
\begin{enumerate}[label=(\alph*)]
    \item What is the distribution of $X$?
    \item Compute $\PP(X > 10)$.
    \item Compute $\PP(X > 15 \mid X > 5)$. Verify the memoryless property.
\end{enumerate}
\end{exercise}

\begin{exercise}
Prove that the geometric distribution is the only discrete distribution on $\N^*$
with the memoryless property.
\textit{Hint:} set $g(n) = \PP(X > n)$ and use the functional equation
$g(m+n) = g(m)g(n)$.
\end{exercise}

\begin{exercise}
Let $N \sim \mathrm{Poisson}(\lambda)$ and $(X_i)_{i \geq 1}$ be i.i.d.\
$\mathrm{Bernoulli}(p)$, independent of $N$. Set $S = \sum_{i=1}^{N} X_i$.
Show that $S \sim \mathrm{Poisson}(\lambda p)$.
\end{exercise}
